\section{System Model and Neural Receiver}
\label{sec:sys_nrx}

%\vspace{-2mm}

\subsection{Uplink Setup}
\label{sec:system_model}
$U$ single-layer \glspl{UE} transmit over the same time--frequency
resources to an $N_{\mathrm{rx}}$-antenna base station. After OFDM demodulation, the
received vector at subcarrier $s$ and OFDM symbol $t$ is
\begin{equation}
\label{eq:mu_mimo}
  \mathbf{y}_{s,t}
  = \sum_{u=1}^{U}
    \mathbf{h}_{s,t,u}\,x_{s,t,u}
    + \mathbf{n}_{s,t},
  \qquad
  \mathbf{n}_{s,t}\sim\mathcal{CN}(\mathbf{0},N_0\mathbf{I}_{N_{\mathrm{rx}}}),
\end{equation}
where $u$ indexes the \glspl{UE},
$x_{s,t,u}\in\mathbb{C}$ is the symbol transmitted by UE $u$,
$\mathbf{h}_{s,t,u}\in\mathbb{C}^{N_{\mathrm{rx}}}$ is the corresponding channel vector, and
$\mathbf{y}_{s,t},\mathbf{n}_{s,t}\in\mathbb{C}^{N_{\mathrm{rx}}}$ are the received
and noise vectors, respectively. Here, $\mathcal{CN}$ denotes a
circularly symmetric complex Gaussian distribution, $N_0$ is the noise
variance per receive component, and $\mathbf{I}_{N_{\mathrm{rx}}}$ is the
$N_{\mathrm{rx}}$-dimensional identity matrix. Payload bits are \gls{LDPC}-encoded
and mapped to \gls{QAM} symbols; known \gls{DMRS} tones provide an
initial \gls{LS} channel estimate. The neural receiver processes the
post-\gls{FFT} data and pilot grid and outputs per-bit \glspl{LLR} for
conventional 5G~NR decoding.

\vspace{-2mm}

\subsection{Neural Receiver Architecture}
\label{sec:architecture}
We adopt this GNN--CNN architecture for MU-\gls{MIMO}
\gls{PUSCH} reception \cite{scamnrx,wiesmayr2024design,neural_rx}. Each resource element has 18
input features. An input CNN with
channel sequence $18{\to}128{\to}128{\to}56$ initializes a state of
dimension $d_s{=}56$. Each of the $N_{\mathrm{it}}{=}8$ unrolled
blocks applies a $56{\to}64{\to}56$ user-mixing \gls{MLP} and a
$114{\to}128{\to}128{\to}56$ Update CNN to exchange information
across users and the time--frequency grid, respectively. A dense
readout of width $128$ produces \glspl{LLR}. Each separable convolution
comprises a depthwise spatial filter and a pointwise channel mixer.
Pruning targets pointwise, aggregation, and non-readout dense kernels.
Depthwise and readout kernels remain dense
(Section~\ref{sec:pruning}). Shared weights accommodate varying user
counts and bandwidths, enabling training on $4$ \glspl{PRB} and
evaluation on a $132$-\gls{PRB} grid.

\vspace{-2mm}

\subsection{Loss and Training Schedule}
\label{sec:loss}
Training minimizes bitwise \gls{BCE} with an auxiliary
channel-estimation loss. At iteration $k$, for coded bit $b\in\{0,1\}$,
let $\hat{L}^{(k)}$ denote the predicted \gls{LLR}, with positive
values favoring $b=1$, and let $\mathbf{h}$ and $\hat{\mathbf{h}}^{(k)}$
denote the true and estimated channel vectors. The objective is
\begin{align}
  \label{eq:total_loss}
    \mathcal{L}
    &= \sum_{k=1}^{N_{\mathrm{it}}}
       \left(
         \mathcal{L}_{\mathrm{BCE}}^{(k)}
         + \lambda\,\mathcal{L}_{\mathrm{MSE}}^{(k)}
       \right), \\
  \label{eq:bce}
    \mathcal{L}_{\mathrm{BCE}}^{(k)}
    &= -\mathbb{E}\!\left[
         b\log\sigma(\hat{L}^{(k)})
         +(1-b)\log\!\bigl(1-\sigma(\hat{L}^{(k)})\bigr)
       \right], \\
  \label{eq:mse}
    \mathcal{L}_{\mathrm{MSE}}^{(k)}
    &= \mathbb{E}\!\bigl[
         \|\hat{\mathbf{h}}^{(k)}-\mathbf{h}\|_{2}^{2}
       \bigr],
  \end{align}
where $\sigma$ is the logistic sigmoid, $\mathbb{E}$ averages over data
bits and channel realizations, and $\lambda>0$ weights the auxiliary
loss. The coefficient $\lambda$ is set to $0.02$ and $0.01$ in
successive training phases. Training uses the open-source
neural-receiver implementation in Sionna \cite{sionna,neural_rx}.
Each mini-batch contains $N_B{=}128$ samples with random active-user
counts, velocities in $[0,56]$\,m/s, and rate-adjusted $E_b/N_0$ drawn
from phase-specific ranges spanning approximately $0$--$15$\,dB. The
\texttt{FP32} baseline is trained with Adam; quantized models
warm-start from that checkpoint.
